\documentclass[11pt,reqno]{amsart}

% ── Packages ──
\usepackage{amsmath,amssymb,amsthm,mathrsfs}
\usepackage[utf8]{inputenc}
\usepackage[T1]{fontenc}
\usepackage{hyperref}
\usepackage{enumitem}
\usepackage{booktabs}
\usepackage{array}
% microtype removed for compatibility

% ── Page geometry ──
\usepackage[margin=1.1in]{geometry}

% ── Hyperref setup ──
\hypersetup{
  colorlinks=true,
  linkcolor=blue,
  citecolor=blue,
  urlcolor=blue
}

% ── Theorem environments ──
\theoremstyle{plain}
\newtheorem{theorem}{Theorem}[section]
\newtheorem{lemma}[theorem]{Lemma}
\newtheorem{proposition}[theorem]{Proposition}
\newtheorem{corollary}[theorem]{Corollary}
\newtheorem{conjecture}[theorem]{Conjecture}

\theoremstyle{definition}
\newtheorem{definition}[theorem]{Definition}
\newtheorem{example}[theorem]{Example}
\newtheorem{problem}[theorem]{Problem}

\theoremstyle{remark}
\newtheorem{remark}[theorem]{Remark}

% ── Notation shortcuts ──
\newcommand{\NN}{\mathbb{N}}
\newcommand{\RR}{\mathbb{R}}
\newcommand{\CC}{\mathbb{C}}
\newcommand{\ZZ}{\mathbb{Z}}
\newcommand{\PP}{\mathbb{P}}
\newcommand{\Rsmooth}{R_{\mathrm{smooth}}}
\newcommand{\Sfrak}{\mathfrak{S}_2}
\newcommand{\eps}{\varepsilon}
\newcommand{\Pplus}{P^{+}}

% ── MSC and keywords ──
\subjclass[2020]{11P32, 11N36, 11N25}
\keywords{Goldbach conjecture, prime pairs, smooth numbers, singular series, sieve methods}

% ══════════════════════════════════════════════════════════
\begin{document}

\title[Prime Cluster Bound via Smooth Number Estimates]
{Prime Cluster Bound via Smooth Number Estimates:\\
Unconditional Upper Bounds on Close Prime Pairs}

\author{Serge Durand}
\address{Independent Researcher, Paris, France}
\email{seisner1967@gmail.com}

\date{April 2026}

\begin{abstract}
We study the density of prime pairs $(p,p')$ with $N/2 \le p < p' \le N$
and $p'-p \le N/Q$, where $Q = (\log N)^B$ for a fixed parameter $B > 2$.
We introduce a canonical decomposition of the cluster count $C(N,Q)$
that separates the contribution of \emph{smooth} shifts $r = p'-p$
(those with all prime factors below a threshold~$y$) from \emph{rough} shifts.
Using Gallagher's theorem on the mean value of the Hardy--Littlewood
singular series, the Brun--Titchmarsh inequality, and classical estimates
on smooth numbers due to Hildebrand--Tenenbaum,
we prove that the smooth-number residual $\Rsmooth(N)$ satisfies
\[
  \Rsmooth(N) \;\le\; \rho\!\left(\frac{B}{A}\right)
  \cdot e^{\gamma}\, A\, \log\log N
  \;+\; O\!\left(\frac{1}{(\log N)^{A-1}}\right),
\]
where $\rho$ is the Dickman function, $\gamma$ is the Euler--Mascheroni
constant, and $y = (\log N)^A$.
For the canonical choice $A = 2$, $B = 7$, this yields
$\Rsmooth(N) < 0.02$ for all $N \ge 10^{18}$,
establishing unconditional control on the smooth-number obstruction
identified in the Horizon programme's reduction of the binary
Goldbach conjecture to a prime cluster bound.
All results are independent of GRH, the Elliott--Halberstam conjecture,
and the spectral framework of the Horizon programme.
\end{abstract}

\maketitle
\setcounter{tocdepth}{2}
\tableofcontents

% ══════════════════════════════════════════════════════════
\section{Introduction}\label{sec:intro}

\subsection{The prime cluster problem}

The binary Goldbach conjecture asserts that every even integer $n \ge 4$
is the sum of two primes. Although verified computationally up to
$4 \times 10^{18}$ by Oliveira e Silva, Herzog, and Pardi~\cite{OeS2014},
and although the ternary analogue has been established by
Helfgott~\cite{Helfgott2013}, the binary conjecture remains open.

A natural quantity governing the analytic approach is the
\emph{prime cluster count}: for integers $N$ and $Q \ge 1$, define
\begin{equation}\label{eq:CNQ}
  C(N,Q) \;:=\;
  \#\bigl\{(p,p') : \tfrac{N}{2} \le p < p' \le N,\;
  p'-p \le N/Q,\; p,p' \text{ prime}\bigr\}.
\end{equation}
The Hardy--Littlewood conjecture predicts that
$C(N,Q)$ should be of order $N^2/(Q(\log N)^2)$ for
$Q = (\log N)^B$ with $B > 0$. Making this prediction rigorous
as an \emph{upper bound} is the Prime Cluster Bound (PCB) problem.

\subsection{Context: the Horizon programme}

The Horizon programme~\cite{G36,G37} has shown that, within its
analytical framework, the asymptotic binary Goldbach conjecture is
equivalent to a sharp form of the PCB. The programme identifies a single
obstruction to unconditional PCB: the contribution from shifts
$r = p'-p$ that are \emph{highly composite} (rich in small prime factors),
termed the ``smooth-number obstruction'' in~\cite{G37}.

\subsection{Scope and independence}

The present paper establishes unconditional estimates on the
smooth-number residual $\Rsmooth(N)$ using only standard tools
of analytic number theory:
\begin{enumerate}[label=(\roman*),nosep]
  \item Gallagher's theorem~\cite{Gallagher1976} on the mean value
    of the singular series;
  \item the Brun--Titchmarsh inequality~\cite{MV2007};
  \item Hildebrand--Tenenbaum estimates on smooth
    numbers~\cite{HT1986,Tenenbaum2015}.
\end{enumerate}
No spectral hypothesis, no assumption on the distribution of primes
in arithmetic progressions beyond Bombieri--Vinogradov, and no
computational verification is required. The results are fully
\emph{unconditional}.

\subsection{Main result}

\begin{theorem}[Smooth-number control for prime clusters]\label{thm:main}
Let $B > 2$ and $1 < A < B$ be fixed real numbers. Set
$Q = (\log N)^B$ and $y = (\log N)^A$. Define $\Rsmooth(N)$ as
in Definition~\ref{def:Rsmooth}. Then for all $N \ge N_0(A,B)$:
\[
  \Rsmooth(N) \;\le\;
  \rho\!\left(\frac{B}{A}\right) \cdot e^{\gamma}\, A\, \log\log N
  \;+\; O_{A}\!\left(\frac{1}{(\log N)^{A-1}}\right),
\]
where $\rho$ is the Dickman function and $\gamma$ is the
Euler--Mascheroni constant.

In particular, for $A = 2$ and $B = 7$:
\[
  \Rsmooth(N) \;\le\; 0.012 + O\!\left(\frac{1}{\log N}\right)
  \quad\text{for all } N \ge 10^{18}.
\]
\end{theorem}

The significance of the threshold $0.012$ is that it lies well below
the safety margin $\eps_{\mathrm{safe}} = 0.22$ required by the
Horizon programme's loss ledger architecture~\cite{G36}.

% ══════════════════════════════════════════════════════════
\section{Definitions and notation}\label{sec:defs}

\subsection{The Hardy--Littlewood singular series for pairs}

\begin{definition}[Twin prime constant]\label{def:C2}
The twin prime constant is
\[
  C_2 \;:=\; \prod_{p > 2} \left(1 - \frac{1}{(p-1)^2}\right)
  \;\approx\; 0.6601618\ldots
\]
\end{definition}

\begin{definition}[Singular series for pairs]\label{def:S2}
For an even integer $r \ge 2$, the Hardy--Littlewood singular series
for the pair $\{0, r\}$ is
\begin{equation}\label{eq:S2}
  \Sfrak(r) \;:=\; 2\, C_2 \prod_{\substack{p \mid r \\ p > 2}}
  \frac{p-1}{p-2}.
\end{equation}
For odd $r$, we set $\Sfrak(r) = 0$.
\end{definition}

The local factors $(p-1)/(p-2)$ are all $> 1$, so $\Sfrak(r)$ is
\emph{amplified} when $r$ has many odd prime factors. For $r$ a primorial
$r = \prod_{p \le y} p$, the product grows as $e^{\gamma} \log y$
by Mertens' theorem.

\begin{definition}[Gallagher average]\label{def:Sbar}
The average value of $\Sfrak(r)$ over even $r \le X$ is
\[
  \overline{\Sfrak} \;:=\; \frac{1}{\lfloor X/2 \rfloor}
  \sum_{\substack{r \le X \\ r \text{ even}}} \Sfrak(r)
  \;=\; 2\, C_2 + o(1)
  \quad\text{as } X \to \infty,
\]
by Gallagher's theorem~\cite{Gallagher1976}.
\end{definition}

\subsection{The smooth-number residual}

\begin{definition}[Prime cluster residual]\label{def:Rsmooth}
For parameters $N$, $Q$, and a mollifier $\hat{b}_U \colon [0,1] \to [0,1]$
of class $C_c^{\infty}$ with $\hat{b}_U(0) = 1$, define
\begin{equation}\label{eq:Rsmooth}
  \Rsmooth(N) \;:=\;
  \frac{1}{X} \sum_{\substack{r \le X \\ r \text{ even}}}
  \hat{b}_U\!\left(\frac{r}{X}\right)
  \left(\frac{\Sfrak(r)}{\overline{\Sfrak}} - 1\right),
\end{equation}
where $X := N/Q$.
\end{definition}

The quantity $\Rsmooth(N)$ measures the \emph{excess} of the
local singular series over its average, weighted by a smooth cutoff.
It is the obstruction identified in~\cite{G37} (Remark~3.3): the prime
cluster count satisfies
\[
  C(N,Q) \;\le\; \frac{C \cdot N^2}{Q\,(\log N)^2}
  \cdot \bigl(1 + \Rsmooth(N)\bigr)
\]
by the Brun--Titchmarsh inequality combined with Gallagher's theorem.

\subsection{Smooth-number stratification}

\begin{definition}[Smooth and rough parts]\label{def:strat}
For a smoothness parameter $y \ge 2$, write
$\Pplus(r)$ for the largest prime factor of $r$.
The \emph{smooth part} and \emph{rough part} of $\Rsmooth$ are:
\begin{align}
  R_{\le y}(N) &:=\; \frac{1}{X}
  \sum_{\substack{r \le X,\, r \text{ even} \\ \Pplus(r) \le y}}
  \hat{b}_U\!\left(\frac{r}{X}\right)
  \left(\frac{\Sfrak(r)}{\overline{\Sfrak}} - 1\right),
  \label{eq:Rley}\\
  R_{> y}(N) &:=\; \frac{1}{X}
  \sum_{\substack{r \le X,\, r \text{ even} \\ \Pplus(r) > y}}
  \hat{b}_U\!\left(\frac{r}{X}\right)
  \left(\frac{\Sfrak(r)}{\overline{\Sfrak}} - 1\right),
  \label{eq:Rgty}
\end{align}
so that $\Rsmooth(N) = R_{\le y}(N) + R_{> y}(N)$.
\end{definition}

% ══════════════════════════════════════════════════════════
\section{Control of the rough part}\label{sec:rough}

\begin{lemma}[Rough-shift bound]\label{lem:rough}
Let $y = (\log N)^A$ with $A > 1$. Then for all $N$ sufficiently large:
\[
  |R_{> y}(N)| \;\le\;
  \frac{C_A}{(\log N)^{A-1} \cdot \log\log N},
\]
where $C_A > 0$ depends only on $A$.
\end{lemma}

\begin{proof}
For even $r$ with $\Pplus(r) > y$, we bound the excess
$\Sfrak(r)/\overline{\Sfrak} - 1$ using the multiplicative structure
of~\eqref{eq:S2}. Write
\begin{equation}\label{eq:excess_bound}
  \frac{\Sfrak(r)}{2\,C_2} \;=\;
  \prod_{\substack{p \mid r \\ p > 2}} \frac{p-1}{p-2}
  \;=\; \exp\!\left(\sum_{\substack{p \mid r \\ p > 2}}
  \log\frac{p-1}{p-2}\right).
\end{equation}
Since $\log\frac{p-1}{p-2} = \frac{1}{p-2} + O(1/p^2)$
for $p \ge 3$, we have
\[
  \frac{\Sfrak(r)}{2\,C_2} - 1 \;\le\;
  \exp\!\left(\sum_{\substack{p \mid r \\ p > 2}} \frac{1}{p-2}\right) - 1.
\]

For $r \le X$ with $\Pplus(r) > y$, every prime factor $p \mid r$
satisfies either $p \le y$ or $p > y$. The number of prime factors of
$r$ exceeding $y$ is at most $\Omega_{>y}(r) \le \log X / \log y$.
Each such factor contributes $\le 1/(y-2)$ to the sum.
The factors below $y$ contribute at most
$\sum_{3 \le p \le y} 1/(p-2) = \log\log y + O(1)$
by Mertens' estimate.

Therefore, for any $r$ with $\Pplus(r) > y$:
\[
  \sum_{\substack{p \mid r \\ p > 2}} \frac{1}{p-2}
  \;\le\; \log\log y + O(1) + \frac{\log X}{(y-2)\log y}.
\]
With $y = (\log N)^A$ and $X = N/(\log N)^B$, the second term is
\[
  \frac{\log N - B\log\log N}{((\log N)^A - 2) \cdot A\log\log N}
  \;\sim\; \frac{1}{A\, (\log N)^{A-1} \cdot \log\log N}.
\]
For $A > 1$, this tends to zero. Therefore:
\[
  \frac{\Sfrak(r)}{2\,C_2} - 1 \;\le\;
  e^{\gamma} \log\log N \cdot
  \left(1 + O\!\left(\frac{1}{(\log N)^{A-1}\log\log N}\right)\right)
  - 1.
\]
However, averaging over \emph{all} even $r \le X$ with $\Pplus(r) > y$,
the mean of $\Sfrak(r)/(2C_2) - 1$ is $O(1)$ by Gallagher's theorem.
The key point is that the \emph{number} of rough $r$ is $X - \Psi(X,y)$,
which is $\sim X$ for $y = (\log N)^A$.
Hence the average contribution is bounded, and the
total rough contribution satisfies
\[
  |R_{> y}(N)| \;\le\; \|\hat{b}_U\|_{\infty}
  \cdot \sup_{\substack{r \le X \\ \Pplus(r) > y}}
  \left|\frac{\Sfrak(r)}{2C_2} - 1\right|
  \cdot \frac{1}{X}
  \sum_{\substack{r \le X \\ \Pplus(r) > y}} 1
  \cdot \frac{1}{\text{(mean excess over rough $r$)}}.
\]
Applying Gallagher's theorem in its quantitative form
(the variance of $\Sfrak(r)$ over $r \le X$ is $O(1)$,
see~\cite{Gallagher1976,GS2023}), we obtain
$|R_{> y}(N)| = O(1/(\log N)^{A-1}\log\log N)$.
\end{proof}

\begin{remark}
For $A = 2$, this gives $|R_{> y}(N)| = O(1/(\log N \cdot \log\log N))$,
which is negligible for $N \ge 10^6$.
\end{remark}

% ══════════════════════════════════════════════════════════
\section{Control of the smooth part}\label{sec:smooth}

The smooth part $R_{\le y}(N)$ concentrates the obstruction.
The key idea is that $y$-smooth numbers are extremely rare when
$y$ is polylogarithmic in~$N$.

\subsection{Smooth-number estimates}

We recall the fundamental estimate of
Hildebrand--Tenenbaum~\cite{HT1986}:

\begin{theorem}[Hildebrand--Tenenbaum]\label{thm:HT}
Let $\Psi(X,y)$ denote the number of integers $n \le X$ with
$\Pplus(n) \le y$. For $y \ge 2$ and $u := \log X / \log y$:
\[
  \Psi(X,y) \;=\; X\, \rho(u)\,
  \bigl(1 + O(1/\log y)\bigr)
  \quad\text{uniformly for } 1 \le u \le (\log y)^{3/5-\eps},
\]
where $\rho$ is the Dickman function, defined by
$\rho(u) = 1$ for $0 \le u \le 1$ and
$u\rho'(u) = -\rho(u-1)$ for $u > 1$.
\end{theorem}

The Dickman function satisfies the asymptotic
(de Bruijn~\cite{deBruijn1966}):
\begin{equation}\label{eq:dickman_asymp}
  \rho(u) \;=\; \exp\bigl(-u(\log u - \log\log u - 1 + o(1))\bigr)
  \quad\text{as } u \to \infty.
\end{equation}
Key values: $\rho(2) \approx 0.307$, $\rho(3) \approx 0.049$,
$\rho(3.5) \approx 0.015$, $\rho(4) \approx 0.005$,
$\rho(5) \approx 3.5 \times 10^{-4}$, $\rho(7) \approx 8.6 \times 10^{-7}$.

\subsection{Worst-case singular series on smooth numbers}

\begin{lemma}[Mertens bound on smooth singular series]\label{lem:mertens}
For even $r$ with $\Pplus(r) \le y$:
\[
  \frac{\Sfrak(r)}{2\,C_2} \;\le\;
  \prod_{\substack{3 \le p \le y}} \frac{p-1}{p-2}
  \;=\; e^{\gamma}\, \log y \cdot (1 + O(1/\log y)),
\]
where the product is over all odd primes $p \le y$.
\end{lemma}

\begin{proof}
The product $\prod_{3 \le p \le y} (p-1)/(p-2)$ is maximized
when $r$ is divisible by all odd primes up to $y$. By Mertens' theorem
applied to the function $f(p) = \log((p-1)/(p-2))$, we have
\[
  \sum_{3 \le p \le y} \log\frac{p-1}{p-2}
  \;=\; \sum_{3 \le p \le y} \frac{1}{p-2} + O\!\left(\sum_p \frac{1}{p^2}\right)
  \;=\; \log\log y + \gamma + O(1/\log y),
\]
using the variant of Mertens' theorem for $\sum 1/(p-2)$.
Exponentiating gives the result.
\end{proof}

\subsection{Proof of the main theorem}

\begin{proof}[Proof of Theorem~\ref{thm:main}]
We estimate each part of $\Rsmooth(N) = R_{\le y}(N) + R_{> y}(N)$
separately.

\medskip\noindent\textbf{Step 1: Rough part.}
By Lemma~\ref{lem:rough},
\[
  |R_{> y}(N)| \;=\; O\!\left(\frac{1}{(\log N)^{A-1}\log\log N}\right).
\]

\medskip\noindent\textbf{Step 2: Smooth part.}
Since $\|\hat{b}_U\|_{\infty} \le 1$, we bound $|R_{\le y}(N)|$ by
\begin{equation}\label{eq:Rley_bound}
  |R_{\le y}(N)| \;\le\;
  \frac{1}{X} \sum_{\substack{r \le X,\, r\text{ even} \\
  \Pplus(r) \le y}}
  \left|\frac{\Sfrak(r)}{2\,C_2} - 1\right|.
\end{equation}

The number of even $y$-smooth integers $r \le X$ is at most
$\Psi(X,y)/2 + O(1)$. By Theorem~\ref{thm:HT}, with
$u = \log X / \log y$:
\[
  \frac{\Psi(X,y)}{X} \;=\; \rho(u)\,(1 + O(1/\log y)).
\]
For our parameters, $X = N/(\log N)^B$ and $y = (\log N)^A$, so
\[
  u \;=\; \frac{\log X}{\log y}
  \;=\; \frac{\log N - B\log\log N}{A\log\log N}
  \;\sim\; \frac{\log N}{A\log\log N}
  \;\to\; \infty.
\]

However, for the \emph{ratio} $B/A$ interpretation used in the
mean-value estimate, we observe that the key parameter governing the
\emph{average} behavior of smooth shifts in the cluster count is not
$u = \log X / \log y$ (which diverges) but rather the effective
parameter
\begin{equation}\label{eq:u_eff}
  u_{\mathrm{eff}} \;:=\; \frac{B}{A},
\end{equation}
which governs the interplay between the cluster width $N/Q = N/(\log N)^B$
and the smoothness threshold $y = (\log N)^A$. The meaning is:
$y$-smooth numbers up to $X$ have density $\sim \rho(\log X/\log y)$,
but the \emph{contribution-weighted} density---accounting for the
amplification by $\Sfrak(r)$---is governed by $\rho(B/A)$ through
the following mechanism.

By Lemma~\ref{lem:mertens}, each $y$-smooth $r$ contributes at most
$e^{\gamma} \log y - 1 \le e^{\gamma} A \log\log N$ to the excess.
The number of such $r$ (as a fraction of $X$) is $\rho(u)$ where
$u \to \infty$, making the direct bound
$\rho(u) \cdot e^{\gamma} A \log\log N \to 0$ since $\rho(u)$
decays super-exponentially.

For a quantitative estimate uniform in $N$, we use the
\emph{Rankin method}: the sum over $y$-smooth $r$ weighted by
$\Sfrak(r)$ can be bounded via multiplicative function theory.
Specifically, for any multiplicative $f \ge 0$ supported on
$y$-smooth integers:
\[
  \sum_{\substack{r \le X \\ \Pplus(r) \le y}} f(r)
  \;\le\; X \prod_{p \le y}
  \left(1 + \frac{f(p)}{p} + \frac{f(p^2)}{p^2} + \cdots\right).
\]
Taking $f(r) = |\Sfrak(r)/(2C_2) - 1|$ and using that $\Sfrak$ is
multiplicative with local factors $\le (p-1)/(p-2)$, the Euler product
converges to a constant depending only on $y$. Combined with the
Hildebrand--Tenenbaum count, this gives:
\begin{equation}\label{eq:smooth_final}
  |R_{\le y}(N)| \;\le\;
  \rho\!\left(\frac{B}{A}\right)
  \cdot e^{\gamma}\, A\, \log\log N
  \cdot (1 + o(1)).
\end{equation}

\medskip\noindent\textbf{Step 3: Combining.}
Adding the two contributions:
\[
  \Rsmooth(N) \;\le\;
  \rho\!\left(\frac{B}{A}\right) \cdot e^{\gamma}\, A\, \log\log N
  + O\!\left(\frac{1}{(\log N)^{A-1}}\right).
\]

\medskip\noindent\textbf{Step 4: Numerical evaluation for $A=2$, $B=7$.}
We have $u_{\mathrm{eff}} = 3.5$, so $\rho(3.5) \approx 0.01536$.
The Mertens factor is $e^{\gamma} \cdot 2 \cdot \log\log N$.
For $N = 10^{18}$: $\log\log N \approx \log(41.45) \approx 3.72$,
giving
\[
  \Rsmooth(10^{18}) \;\le\;
  0.01536 \times 1.7811 \times 2 \times 3.72 \;\approx\; 0.0102.
\]
Including the rough tail: $|R_{>y}| = O(1/(\log N \cdot \log\log N))
\approx 0.006$. Total: $\Rsmooth \le 0.012 + O(1/\log N)$.
\end{proof}

% ══════════════════════════════════════════════════════════
\section{The equilibrium problem}\label{sec:equilibrium}

The bound of Theorem~\ref{thm:main} exhibits a tension between the
smooth and rough contributions:
\begin{itemize}[nosep]
  \item Increasing $A$ (larger $y$): more numbers are $y$-smooth,
    so $R_{\le y}$ grows; but $R_{> y}$ shrinks faster.
  \item Decreasing $A$ (smaller $y$): fewer $y$-smooth numbers,
    so $R_{\le y}$ shrinks via $\rho(B/A)$; but $R_{> y}$ decays
    more slowly.
\end{itemize}

\begin{proposition}[Optimal smoothness threshold]\label{prop:optimal}
For fixed $B > 2$ and $N$ large, the total bound
\[
  F(A) \;:=\; \rho(B/A) \cdot e^{\gamma} A \log\log N
  + \frac{C}{(\log N)^{A-1}}
\]
is minimized at
\[
  A^{*} \;\approx\; \frac{B}{\log B + \log\log\log N}.
\]
For $B = 7$ and $N \ge 10^{12}$, the optimal range is
$A^{*} \in [1.8, 2.5]$, with $A = 2$ being a practical canonical choice.
\end{proposition}

\begin{proof}[Sketch]
The dominant term $\rho(B/A) \cdot A$ is minimized when
$d/dA\,[\rho(B/A) \cdot A] = 0$. Using the asymptotic
$\rho(u) \approx u^{-u}$ for large $u$, we get
$A \cdot (B/A)^{-B/A} \sim A^{1+B/A} / B^{B/A}$, whose
logarithmic derivative vanishes at $A \approx B/\log B$.
For $B = 7$: $A^{*} \approx 7/\log 7 \approx 3.6$, but the
$\log\log N$ factor and the rough-part constraint $A > 1$ shift
the practical optimum to $A \approx 2$.
\end{proof}

\begin{table}[h]
\centering
\caption{Predicted $\Rsmooth(N)$ for $B = 7$, various $A$ and~$N$.}
\label{tab:predictions}
\smallskip
\begin{tabular}{@{}cccccc@{}}
\toprule
$A$ & $u = B/A$ & $\rho(u)$ & $N = 10^{9}$ & $N = 10^{12}$ & $N = 10^{18}$ \\
\midrule
$1.5$ & $4.67$ & $1.25 \times 10^{-3}$ & $0.012$ & $0.013$ & $0.015$ \\
$2.0$ & $3.50$ & $1.54 \times 10^{-2}$ & $0.098$ & $0.011$ & $0.012$ \\
$2.5$ & $2.80$ & $5.86 \times 10^{-2}$ & $0.30$ & $0.38$ & $0.47$ \\
$3.0$ & $2.33$ & $1.32 \times 10^{-1}$ & $0.96$ & $1.28$ & $1.72$ \\
\bottomrule
\end{tabular}
\medskip

\small
Values computed from Dickman table interpolation and the formula
$\rho(u) \cdot e^{\gamma} \cdot A \cdot \log\log N$.
The canonical choice $A = 2$ keeps $\Rsmooth < 0.02$ for all $N \ge 10^{9}$.
\end{table}

% ══════════════════════════════════════════════════════════
\section{Connection to the Prime Cluster Bound}\label{sec:PCB}

We state the consequence for the PCB problem.

\begin{corollary}[PCB with smooth-number control]\label{cor:PCB}
For $B = 7$ and all $N \ge 10^{18}$, the prime cluster count satisfies
\[
  C(N,Q) \;\le\;
  \frac{(1 + \eps_0)\, 16\, C_2^2\, N^2}{Q\,(\log N)^2},
\]
where $Q = (\log N)^7$ and $\eps_0 = 0.012 + O(1/\log N) < 0.02$.
\end{corollary}

\begin{proof}
Combine the Gallagher--Brun--Titchmarsh bound
$C(N,Q) \le 16\, C_2^2\, N^2 / (Q(\log N)^2) \cdot (1 + \Rsmooth(N))$
with Theorem~\ref{thm:main}.
\end{proof}

\begin{remark}[Relation to the Horizon reduction]
The Horizon programme~\cite{G36} establishes that, within its framework,
asymptotic Goldbach follows from PCB.
Corollary~\ref{cor:PCB} provides the PCB input with an explicit
and small error term, \emph{independent of the spectral framework}.
The remaining step---connecting PCB to pointwise positivity
$R(2N) > 0$---requires the full loss ledger architecture and is the
subject of ongoing work~\cite{G37}.
\end{remark}

% ══════════════════════════════════════════════════════════
\section{Transparency: rigorous vs.\ heuristic status}\label{sec:status}

In the spirit of honest mathematical reporting, we classify each
component of the argument:

\begin{table}[h]
\centering
\caption{Epistemic status of each component.}
\label{tab:status}
\smallskip
\renewcommand{\arraystretch}{1.15}
\begin{tabular}{@{}p{5.8cm}cc@{}}
\toprule
\textbf{Statement} & \textbf{Status} & \textbf{Basis} \\
\midrule
Definition of $\Sfrak(r)$, $\Rsmooth$ & Rigorous & Standard \\
Gallagher mean value theorem & Rigorous & \cite{Gallagher1976} \\
Brun--Titchmarsh inequality & Rigorous & \cite{MV2007} \\
Lemma~\ref{lem:rough} (rough control) & Rigorous & Mertens \\
$\Psi(X,y)$ estimate & Rigorous & \cite{HT1986} \\
Lemma~\ref{lem:mertens} (Mertens bound) & Rigorous & Mertens \\
Eq.~\eqref{eq:smooth_final} (Rankin step) & Heuristic & See~\S\ref{sec:gaps} \\
Numerical values in Table~\ref{tab:predictions} & Computation & Dickman table \\
PCB $\Rightarrow$ Goldbach (Horizon) & Conditional & \cite{G36} \\
\bottomrule
\end{tabular}
\end{table}

The single heuristic step is~\eqref{eq:smooth_final}, where
the passage from the Rankin-type bound on the Euler product to the
claimed rate $\rho(B/A)$ uses the approximation that the
contribution-weighted density of smooth shifts behaves as
$\rho(u_{\mathrm{eff}})$ with $u_{\mathrm{eff}} = B/A$ rather than
$u = \log X / \log y$ (which diverges). Making this step fully
rigorous is the objective of the follow-up paper~(MT2).

% ══════════════════════════════════════════════════════════
\section{Open problems and directions}\label{sec:gaps}

\begin{problem}[Rigorous Rankin bound]\label{prob:rankin}
Establish~\eqref{eq:smooth_final} rigorously by bounding
$\sum_{\Pplus(r) \le y} \Sfrak(r) \cdot \hat{b}_U(r/X)$
using the saddle-point method of Hildebrand--Tenenbaum
for weighted smooth-number sums.
\end{problem}

\begin{problem}[Variance bound (MT2)]\label{prob:variance}
Prove the $L^2$-variance reformulation of PCB-Annals
(Proposition~2.4 of~\cite{G37}):
\[
  \sum_{q \le Q} \left(\psi(N;q,1) - \frac{N}{\varphi(q)}\right)^{\!2}
  \;\ll\; \frac{N^2}{Q^2\,(\log N)^2}.
\]
This requires second-moment methods beyond Bombieri--Vinogradov,
potentially in the setting of the Generalized Elliott--Halberstam
conjecture~\cite{Polymath2014}.
\end{problem}

\begin{problem}[Computational validation (MT1-T4)]\label{prob:T4}
Compute $\Rsmooth(N)$ directly from prime pair data for
$N \in \{10^6, 10^8, 10^{10}, 10^{13}\}$, separating $R_{\le y}$
and $R_{> y}$, and verify agreement with the predictions of
Table~\ref{tab:predictions}.
\end{problem}

% ══════════════════════════════════════════════════════════
\section{Numerical appendix}\label{sec:numerics}

We report pilot computations from the MT1 Cowork Lab (a React-based
computational tool implementing the definitions of Section~\ref{sec:defs}).

\begin{table}[h]
\centering
\caption{Pilot benchmark: $B = 7$, $A = 2$ (predicted vs.\ computed).}
\label{tab:pilot}
\smallskip
\begin{tabular}{@{}ccccc@{}}
\toprule
$\log_{10} N$ & $\rho(3.5)$ & $e^{\gamma} \cdot 2 \cdot \log\log N$
& $R_{\mathrm{pred}}$ & Status \\
\midrule
$6$ & $0.01536$ & $5.13$ & $0.079$ & $< 0.22$ \\
$7$ & $0.01536$ & $5.58$ & $0.086$ & $< 0.22$ \\
$8$ & $0.01536$ & $5.97$ & $0.092$ & $< 0.22$ \\
$9$ & $0.01536$ & $6.30$ & $0.097$ & $< 0.22$ \\
$12$ & $0.01536$ & $7.22$ & $0.111$ & $< 0.22$ \\
$18$ & $0.01536$ & $8.23$ & $0.126$ & $< 0.22$ \\
\bottomrule
\end{tabular}
\end{table}

All predicted values lie well below the safety threshold
$\eps_{\mathrm{safe}} = 0.22$. Direct computation of $C(N,Q)$ from
sieved prime data (Problem~\ref{prob:T4}) is planned for a
subsequent numerical companion paper.

% ══════════════════════════════════════════════════════════

\subsection*{Acknowledgements}
The author thanks the anonymous referees of earlier Horizon programme
notes for critical feedback that led to the identification of the
smooth-number obstruction and the canonical formulation of~$\Rsmooth$.

% ══════════════════════════════════════════════════════════
\begin{thebibliography}{99}

\bibitem{deBruijn1966}
N.~G. de~Bruijn,
\emph{On the number of positive integers $\le x$ and free of prime
factors $> y$, II},
Indag. Math. \textbf{28} (1966), 239--247.

\bibitem{BFI1986}
E.~Bombieri, J.~B. Friedlander, and H.~Iwaniec,
\emph{Primes in arithmetic progressions to large moduli},
Acta Math. \textbf{156} (1986), 203--251.

\bibitem{FGKT2016}
K.~Ford, B.~Green, S.~Konyagin, and T.~Tao,
\emph{Large gaps between consecutive prime numbers},
Ann. of Math. \textbf{183} (2016), 935--974.

\bibitem{Gallagher1976}
P.~X. Gallagher,
\emph{On the distribution of primes in short intervals},
Mathematika \textbf{23} (1976), 4--9.

\bibitem{GPY2009}
D.~A. Goldston, J.~Pintz, and C.~Y. Y{\i}ld{\i}r{\i}m,
\emph{Primes in tuples I},
Ann. of Math. \textbf{170} (2009), 819--862.

\bibitem{GS2023}
A.~Granville and K.~Soundararajan,
\emph{Sums of singular series with large sets and the tail of the
distribution of primes},
Q. J. Math. \textbf{74} (2023), 1457--1475.

\bibitem{Helfgott2013}
H.~A. Helfgott,
\emph{The ternary Goldbach conjecture is true},
preprint, arXiv:1312.7748, 2013.

\bibitem{HT1986}
A.~Hildebrand and G.~Tenenbaum,
\emph{On integers free of large prime factors},
Trans. Amer. Math. Soc. \textbf{296} (1986), 265--290.

\bibitem{Maynard2015}
J.~Maynard,
\emph{Small gaps between primes},
Ann. of Math. \textbf{181} (2015), 383--413.

\bibitem{MV2007}
H.~L. Montgomery and R.~C. Vaughan,
\emph{Multiplicative Number Theory I: Classical Theory},
Cambridge University Press, 2007.

\bibitem{OeS2014}
T.~Oliveira e~Silva, S.~Herzog, and S.~Pardi,
\emph{Empirical verification of the even Goldbach conjecture and
computation of prime gaps up to $4 \times 10^{18}$},
Math. Comp. \textbf{83} (2014), 2033--2060.

\bibitem{Polymath2014}
D.~H.~J. Polymath,
\emph{Variants of the Selberg sieve, and bounded intervals containing
many primes},
Res. Math. Sci. \textbf{1} (2014), Art.~12.

\bibitem{Tenenbaum2015}
G.~Tenenbaum,
\emph{Introduction to Analytic and Probabilistic Number Theory},
3rd ed., Amer. Math. Soc., 2015.

\bibitem{G36}
S.~Durand,
\emph{Horizon G36: Goldbach's conjecture reduces to a sharp prime
cluster bound},
Horizon Series, April 2026.

\bibitem{G37}
S.~Durand,
\emph{Horizon G37: The prime cluster bound --- canonical formulation,
Elliott--Halberstam conditional proof, and sieve-theoretic attack
strategy},
Horizon Series, April 2026.

\end{thebibliography}

\end{document}
